%% ===== Abstract =====
\begin{abstract}
Large language models (LLMs) are now used to translate code between languages and to migrate it across versions, libraries and frameworks. Whether the \new{migrated code} behaves like the original is always decided by some oracle, and the choice of oracle can \new{change} the reported result. We present a systematic mapping study of how behavioral equivalence is evaluated in this \new{line of research}, covering the main proceedings of ICSE, FSE and ASE from 2024 to 2026. From 4,451 papers we screened \new{431} candidates and included 51 (49 in the main corpus, 2 on cross-language equivalence detection), each \new{read in full}. Three findings stand out. First, compilation (67\%) and existing tests (57\%) dominate, while formal verification (4\%) and LLM judges (6\%) are rare; 9 papers use no execution-based oracle at all. Second, the oracle follows the context: no legacy migration uses existing tests, library migration relies on human judgment (9 of 12 papers), and LLM-generated tests and LLM judges appear only in 2026. Third, \new{33 papers report two or more oracles on the same outputs; among the 18 that report both a compile rate and a behavioral pass rate, compile rates exceed behavioral pass rates by up to 98 percentage points. Nine} papers state no limitation of their own oracle. We contribute a coded map of oracle practice, \new{a catalog of how oracles disagree on the same outputs, a classification of the limitations studies report about their own oracles, six gaps between evaluation practice and real migration, and a checklist for reporting oracle provenance.}
\end{abstract}
